\begin{textsample}{POS Dim 3 – human – Score 8.00 – mg/mg\_ef\_linguagens\_arte\_1\_p7.txt}  \label{ex:f3_pos_002}
Aliada a ela, e tendo como fundamento a BNCC, que desde 2017 se tornou lei, o presente currículo tem intenção de concretizar a \textbf{primeira} iniciativa da história do nosso país de construção de um \textbf{documento} normativo referência para os Projetos Políticos Pedagógicos das escolas. \textbf{Cabe} lembrar que, apesar da iniciativa inédita, no nível \textbf{nacional}, Minas Gerais já contava há mais de dez anos com uma \textbf{proposta} \textbf{curricular} estadual, os Conteúdos Básicos \textbf{Comuns} ( CBCs ), que, após revisão em 2014, passou a ser denominada Currículo Básico Comum ( CBC ), e direcionaram o trabalho da educação pública estadual em todos os componentes \textbf{curriculares}, tendo por sua vez os Parâmetros \textbf{Curriculares} Nacionais ( PCNs ) como ponto de partida.
E no caso do componente \textbf{curricular} Arte não era diferente.

% matched lemmas: caber, comum, curricular, documento, nacional, primeiro, proposta
\end{textsample}
